\chapter{The Heston Stochastic Volatility Data Simulator}
\label{chap:heston}

This chapter provides a self-contained technical exposition of the \texttt{HestonSimulator} class, a data-generation object that produces trajectories of financial log-prices under the Heston (1993) stochastic volatility model and computes the conditional probability density function (PDF) of future log-prices using the Fractional Fast Fourier Transform (frFFT). Every mathematical concept required to understand the implementation is introduced from first principles.

% ─────────────────────────────────────────────────────────────────────────────
\section{The Heston Stochastic Volatility Model}
\label{sec:heston_model}

\subsection{Limitations of Constant-Volatility Models}

The classical Black--Scholes model assumes that the instantaneous volatility of an asset price is a deterministic constant $\sigma$. This assumption leads to a log-normal distribution for the terminal asset price, which is analytically tractable but empirically inadequate: financial markets exhibit a so-called \textit{volatility smile}, whereby options at different strikes imply different values of $\sigma$ when inverted through the Black--Scholes formula. This observation alone establishes that volatility cannot be constant.

\subsection{The Heston System of SDEs}

To remedy this shortcoming, \cite{heston1993} introduced a two-factor model in which both the log-price $X_t = \ln S_t$ and the instantaneous variance $v_t$ are stochastic. The system of stochastic differential equations (SDEs) in the \texttt{HestonSimulator} is written in \textit{log-price form}:

\begin{align}
    dX_t &= \left(\mu - \tfrac{1}{2}v_t\right)dt + \sqrt{v_t}\, dW_t^X \label{eq:logprice_sde}\\
    dv_t &= \kappa(\theta - v_t)\,dt + \sigma_v\sqrt{v_t}\, dW_t^v \label{eq:var_sde}\\
    \langle dW_t^X,\, dW_t^v \rangle &= \rho\, dt \label{eq:correlation}
\end{align}

Each term carries a precise financial interpretation.

\paragraph{Log-price SDE (Eq.\,\ref{eq:logprice_sde}).} Applying It\^{o}'s lemma to $X_t = \ln S_t$ under the assumption that $S_t$ follows a geometric Brownian motion with stochastic diffusion coefficient $\sqrt{v_t}$ yields the drift $\mu - v_t/2$. The term $\mu$ is the risk-neutral drift (typically the risk-free rate), while the $-v_t/2$ correction, called the \textit{It\^{o} correction}, arises from the quadratic variation of $X_t$ and ensures that $\mathbb{E}[S_t] = S_0 e^{\mu t}$.

\paragraph{Variance SDE (Eq.\,\ref{eq:var_sde}).} The variance $v_t$ evolves according to a \textit{Cox--Ingersoll--Ross} (CIR) process \cite{cox1985}. The three parameters have well-defined roles: $\kappa > 0$ is the \textit{mean-reversion speed}, controlling how quickly $v_t$ is pulled back toward its long-run equilibrium; $\theta > 0$ is the \textit{long-run variance}, the level to which $v_t$ reverts; and $\sigma_v > 0$ is the \textit{volatility of volatility} (vol-of-vol), governing the magnitude of random fluctuations in $v_t$.

\paragraph{Correlation (Eq.\,\ref{eq:correlation}).} The two Brownian motions $W^X$ and $W^v$ are correlated with instantaneous correlation $\rho \in (-1, 1)$. For equities, $\rho$ is typically negative: a negative shock to the asset price tends to be accompanied by an increase in variance, a well-documented empirical phenomenon known as the \textit{leverage effect}.

\subsection{The Feller Condition}

A key requirement for the CIR process is that $v_t \geq 0$ at all times. \cite{feller1951} proved that the boundary $v_t = 0$ is never reached (i.e., $v_t > 0$ almost surely) if and only if

\begin{equation}
    2\kappa\theta > \sigma_v^2.
    \label{eq:feller}
\end{equation}

When condition (\ref{eq:feller}) is violated, the process touches zero with positive probability. In the simulator, this boundary is handled numerically by the \textit{full-truncation} scheme described in Section~\ref{sec:discretisation}.

% ─────────────────────────────────────────────────────────────────────────────
\section{Correlated Brownian Increments}
\label{sec:brownian}

To simulate the correlated pair $(dW_t^X, dW_t^v)$ over a discrete time step $\Delta t$, one exploits the fact that any bivariate Gaussian vector with correlation $\rho$ can be constructed from two independent standard normals $Z_1, Z_2 \sim \mathcal{N}(0,1)$ via the \textit{Cholesky decomposition}. Specifically, the increments

\begin{align}
    \Delta W^X &= \sqrt{\Delta t}\, Z_1 \label{eq:dWX}\\
    \Delta W^v &= \sqrt{\Delta t}\left(\rho\, Z_1 + \sqrt{1-\rho^2}\, Z_2\right) \label{eq:dWv}
\end{align}

satisfy $\mathbb{E}[\Delta W^X] = \mathbb{E}[\Delta W^v] = 0$, $\text{Var}(\Delta W^X) = \text{Var}(\Delta W^v) = \Delta t$, and $\text{Cov}(\Delta W^X, \Delta W^v) = \rho\,\Delta t$, which matches the covariance structure prescribed by Eq.~(\ref{eq:correlation}). This construction is implemented directly in \texttt{get\_paths()} (lines 225--228 of the source):

\begin{equation}
    \texttt{dW\_v} = \sqrt{\Delta t}\bigl(\rho\cdot Z_1 + \sqrt{1-\rho^2}\cdot Z_2\bigr).
\end{equation}

The decomposition is exact regardless of the sign or magnitude of $\rho$, provided $|\rho| < 1$.

% ─────────────────────────────────────────────────────────────────────────────
\section{Numerical Discretisation}
\label{sec:discretisation}

Equations (\ref{eq:logprice_sde})--(\ref{eq:var_sde}) do not admit closed-form solutions in general, so they must be integrated numerically. The simulator supports two schemes.

\subsection{Euler--Maruyama for the Log-Price}

The \textit{Euler--Maruyama} scheme \cite{kloeden1992} is the SDE analogue of Euler's method for ODEs. For a generic It\^{o} SDE $dY_t = a(Y_t)\,dt + b(Y_t)\,dW_t$, the scheme reads

\begin{equation}
    Y_{n+1} = Y_n + a(Y_n)\,\Delta t + b(Y_n)\,\Delta W_n.
    \label{eq:euler}
\end{equation}

Applied to the log-price SDE (\ref{eq:logprice_sde}), with $\hat{v}_n = \max(v_n, 0)$ (the truncated variance):

\begin{equation}
    X_{n+1} = X_n + \left(\mu - \tfrac{1}{2}\hat{v}_n\right)\Delta t + \sqrt{\hat{v}_n}\,\Delta W_n^X.
    \label{eq:euler_logprice}
\end{equation}

The Euler--Maruyama scheme has strong convergence order $1/2$ and weak convergence order $1$ for SDEs with globally Lipschitz coefficients.

\subsection{Milstein Scheme with Full Truncation for the Variance}

The \textit{Milstein scheme} \cite{milstein1974} improves upon Euler--Maruyama by including an additional correction term derived from It\^{o}'s formula applied to the diffusion coefficient. For an SDE with diffusion $b(Y) = \sigma_v\sqrt{Y}$, the correction is

\begin{equation}
    \frac{1}{2}b(Y_n)\frac{\partial b}{\partial Y}(Y_n)\bigl((\Delta W_n)^2 - \Delta t\bigr) = \frac{\sigma_v^2}{4}\bigl((\Delta W_n^v)^2 - \Delta t\bigr).
    \label{eq:milstein_correction}
\end{equation}

This term captures the curvature of the diffusion coefficient and raises the strong convergence order from $1/2$ to $1$.

\paragraph{Full truncation.} Because the CIR process can reach zero when the Feller condition fails (or nearly so), a naive Milstein step may produce $v_{n+1} < 0$, which would make $\sqrt{v_{n+1}}$ undefined at the next step. The \textit{full-truncation} strategy \cite{lord2010} addresses this in two stages:

\begin{enumerate}
    \item Before computing the drift and diffusion, replace $v_n$ with $\hat{v}_n = \max(v_n, 0)$ inside all coefficient functions.
    \item After the update step, apply a second truncation: $v_{n+1} \leftarrow \max(v_{n+1}, 0)$.
\end{enumerate}

The complete full-truncation Milstein step for the variance is therefore

\begin{equation}
    v_{n+1} = \max\!\Bigl(v_n + \kappa(\theta - \hat{v}_n)\Delta t + \sigma_v\sqrt{\hat{v}_n}\,\Delta W_n^v + \tfrac{1}{4}\sigma_v^2\bigl((\Delta W_n^v)^2 - \Delta t\bigr),\; 0\Bigr).
    \label{eq:milstein_full}
\end{equation}

This approach is implemented at lines 242--253 of the source, where the \texttt{scheme} flag selects between Milstein (default) and pure Euler--Maruyama (which omits the last correction term).

% ─────────────────────────────────────────────────────────────────────────────
\section{The Characteristic Function}
\label{sec:cf}

\subsection{Definition and Properties}

The \textit{characteristic function} (CF) of a real-valued random variable $X$ is defined as

\begin{equation}
    \varphi(u) = \mathbb{E}\!\left[e^{iuX}\right], \quad u \in \mathbb{R},
    \label{eq:cf_def}
\end{equation}

where $i = \sqrt{-1}$. The CF is the Fourier transform of the PDF and encodes all information about the distribution of $X$. Its fundamental importance stems from the fact that many stochastic models—including the Heston model—admit a closed-form CF even when the PDF itself does not exist in closed form. If $\varphi$ is known analytically, the PDF can be recovered numerically via Fourier inversion (Section~\ref{sec:fourier_inv}).

\subsection{The Heston Characteristic Function}

For the Heston model, \cite{heston1993} derived the conditional CF of the log-price $X_{t+\tau}$ given the state $(X_t, v_t)$ at time $t$. The CF takes the affine-exponential form

\begin{equation}
    \varphi(u;\tau) = \exp\!\Bigl(iuX_0 + iu\mu\tau + A(u,\tau) + B(u,\tau)\,v_0\Bigr),
    \label{eq:cf_heston}
\end{equation}

where $X_0 = X_t$ is the current log-price, $v_0 = v_t$ is the current variance, and $\tau$ is the forecast horizon. The functions $A$ and $B$ depend on the model parameters through an intermediate quantity $d$:

\begin{align}
    d &= \sqrt{(\kappa - \rho\sigma_v\, iu)^2 + \sigma_v^2(u^2 + iu)}, \quad \text{Re}(d) \geq 0, \label{eq:d_heston}\\
    \tilde{g} &= \frac{\kappa - \rho\sigma_v\, iu - d}{\kappa - \rho\sigma_v\, iu + d}, \label{eq:g_tilde}\\
    A(u,\tau) &= \frac{\kappa\theta}{\sigma_v^2}\left[(\kappa - \rho\sigma_v\, iu - d)\tau - 2\ln\!\frac{1 - \tilde{g}\,e^{-d\tau}}{1 - \tilde{g}}\right], \label{eq:A_heston}\\
    B(u,\tau) &= \frac{v_0}{\sigma_v^2}(\kappa - \rho\sigma_v\, iu - d)\frac{1 - e^{-d\tau}}{1 - \tilde{g}\,e^{-d\tau}}. \label{eq:B_heston}
\end{align}

\subsection{The ``Little Trap'' Formulation}

The original Heston CF involves a complex-valued square root and a complex logarithm. Because the complex logarithm is multi-valued, numerical implementations must track the correct branch, which can lead to discontinuities in the CF as $u$ varies—a phenomenon known as the \textit{branch-cut problem}. This instability causes spurious oscillations in the numerically inverted PDF.

\cite{albrecher2007} (the ``Little Trap'' paper) showed that a simple reformulation of Eq.~(\ref{eq:g_tilde}) eliminates the branch-cut issue entirely. The key insight is to define $\tilde{g}$ as the ratio $(\xi - d)/(\xi + d)$ where $\xi = \kappa - \rho\sigma_v\, iu$, and to enforce $\text{Re}(d) \geq 0$ in Eq.~(\ref{eq:d_heston}) by choosing the branch of the square root with non-negative real part. Under this convention, the logarithm in Eq.~(\ref{eq:A_heston}) remains on a consistent branch for all $u$, and the CF is smooth and numerically stable. This is implemented in \texttt{\_heston\_cf()} at lines 318--323:

\begin{equation}
    d \;\leftarrow\; \begin{cases} d & \text{if } \text{Re}(d) \geq 0, \\ -d & \text{otherwise.} \end{cases}
\end{equation}

The code evaluates Eqs.~(\ref{eq:d_heston})--(\ref{eq:B_heston}) in vectorised form using NumPy broadcasting: parameters of shape $(J, 1)$ are combined with the frequency grid of shape $(1, N_\text{fft})$ to produce the CF array of shape $(J, N_\text{fft})$ in a single pass over all $J$ trajectories and all $N_\text{fft}$ frequency points.

% ─────────────────────────────────────────────────────────────────────────────
\section{Fourier Inversion of the Characteristic Function}
\label{sec:fourier_inv}

\subsection{The Gil-Pelaez Inversion Formula}

Given the CF $\varphi(u)$ of a continuous random variable $X$ with PDF $f(x)$, the \textit{Gil-Pelaez inversion theorem} \cite{gilpelaez1951} states that

\begin{equation}
    f(x) = \frac{1}{\pi}\int_0^\infty \text{Re}\!\left[\varphi(u)\,e^{-iux}\right]du.
    \label{eq:gilpelaez}
\end{equation}

This follows from the standard Fourier inversion theorem and the fact that $\varphi(-u) = \overline{\varphi(u)}$ for real-valued distributions, which halves the integration range. To evaluate Eq.~(\ref{eq:gilpelaez}) numerically, the integral is truncated at a finite upper limit $U = (N_\text{fft}-1)\eta$ and discretised on a uniform grid $u_j = j\eta$, $j = 0, 1, \ldots, N_\text{fft}-1$, with frequency step $\eta > 0$:

\begin{equation}
    f(x) \approx \frac{\eta}{\pi}\,\text{Re}\!\left[\sum_{j=0}^{N_\text{fft}-1} w_j\,\varphi(j\eta)\,e^{-ij\eta x}\right],
    \label{eq:gilpelaez_discrete}
\end{equation}

where $w_j$ are numerical quadrature weights (Section~\ref{sec:simpson}).

\subsection{Simpson's Rule Quadrature Weights}
\label{sec:simpson}

The choice of quadrature weights $w_j$ in Eq.~(\ref{eq:gilpelaez_discrete}) determines the accuracy of the numerical integral. Simple rectangular (Euler) weights give first-order accuracy. The implementation instead uses \textit{Simpson's rule}, which achieves fourth-order accuracy by approximating the integrand as a piecewise quadratic polynomial. For $N_\text{fft}$ equally spaced points with step $\eta$, Simpson's rule assigns weights

\begin{equation}
    w_j = \frac{1}{3}\begin{cases} 1 & j = 0 \text{ or } j = N_\text{fft}-1, \\ 4 & j \text{ odd}, \\ 2 & j \text{ even}, \; j \neq 0, N_\text{fft}-1, \end{cases}
    \label{eq:simpson}
\end{equation}

which is exactly what lines 498--501 of the source code implement by setting \texttt{w[1:-1:2] = 4.0} and \texttt{w[2:-2:2] = 2.0} before dividing by 3.

% ─────────────────────────────────────────────────────────────────────────────
\section{From Standard FFT to Fractional FFT}
\label{sec:fft}

\subsection{The Discrete Fourier Transform}

The \textit{Discrete Fourier Transform} (DFT) of a sequence $\{x_j\}_{j=0}^{N-1}$ is defined as

\begin{equation}
    \hat{x}_k = \sum_{j=0}^{N-1} x_j\, e^{-i2\pi jk/N}, \quad k = 0, 1, \ldots, N-1.
    \label{eq:dft}
\end{equation}

The \textit{Fast Fourier Transform} (FFT) computes Eq.~(\ref{eq:dft}) in $\mathcal{O}(N\log N)$ operations rather than the $\mathcal{O}(N^2)$ required by direct summation. To apply the FFT to the inversion formula (\ref{eq:gilpelaez_discrete}), one evaluates Eq.~(\ref{eq:gilpelaez_discrete}) simultaneously at the grid of log-price values $x_k = b + k\lambda$, $k = 0, \ldots, N_\text{fft}-1$, for some lower bound $b$ and spacing $\lambda$. Inserting this into Eq.~(\ref{eq:gilpelaez_discrete}):

\begin{align}
    f(x_k) &\approx \frac{\eta}{\pi}\,\text{Re}\!\left[e^{-ij\eta\, b}\sum_{j=0}^{N_\text{fft}-1} w_j\,\varphi(j\eta)\,e^{-i2\pi\cdot\frac{\eta\lambda}{2\pi}\cdot jk}\right].
    \label{eq:fft_coupl}
\end{align}

Comparing with Eq.~(\ref{eq:dft}), the sum is a DFT if and only if the exponent is $e^{-i2\pi jk/N_\text{fft}}$, which requires

\begin{equation}
    \frac{\eta\lambda}{2\pi} = \frac{1}{N_\text{fft}} \;\Longrightarrow\; \lambda = \frac{2\pi}{N_\text{fft}\,\eta}.
    \label{eq:std_fft_constraint}
\end{equation}

This \textit{coupling constraint} links the frequency step $\eta$ and the spatial step $\lambda$: making the log-price grid finer (smaller $\lambda$) forces a larger $\eta$, which reduces the effective integration range in frequency and may introduce aliasing or truncation errors. This rigid coupling is the central limitation of the standard FFT approach to Fourier inversion.

\subsection{The Fractional DFT and Bluestein's Algorithm}
\label{sec:frfft}

The \textit{Fractional DFT} (frDFT) generalises Eq.~(\ref{eq:dft}) by replacing $1/N$ with an arbitrary real parameter $\alpha$:

\begin{equation}
    y_k = \sum_{j=0}^{N-1} x_j\, e^{-i2\pi\alpha jk}, \quad k = 0, \ldots, N-1.
    \label{eq:frdft}
\end{equation}

Comparing with Eq.~(\ref{eq:fft_coupl}), choosing $\alpha = \eta\lambda/(2\pi)$ evaluates the inversion integral at the spatial grid $\{x_k\}$ without any constraint between $\eta$ and $\lambda$. The user is therefore free to choose $\eta$ and $\lambda$ independently, completely decoupling frequency resolution from spatial resolution.

\paragraph{The Bluestein identity.} Computing Eq.~(\ref{eq:frdft}) naively requires $\mathcal{O}(N^2)$ operations. \cite{bluestein1970} observed that the product $jk$ can be written as

\begin{equation}
    jk = \frac{j^2 + k^2 - (j-k)^2}{2},
    \label{eq:bluestein_id}
\end{equation}

which transforms Eq.~(\ref{eq:frdft}) into

\begin{equation}
    y_k = e^{-i\pi\alpha k^2}\sum_{j=0}^{N-1}\!\left[x_j e^{-i\pi\alpha j^2}\right]\cdot e^{i\pi\alpha(j-k)^2}.
    \label{eq:bluestein_conv}
\end{equation}

The sum over $j$ in Eq.~(\ref{eq:bluestein_conv}) is a \textit{linear convolution} of the sequence $a_j = x_j e^{-i\pi\alpha j^2}$ with the symmetric \textit{chirp kernel} $h_m = e^{i\pi\alpha m^2}$, evaluated at lags $m = j - k$. Linear convolution of length-$N$ sequences can be computed in $\mathcal{O}(N\log N)$ via the convolution theorem:

\begin{equation}
    \text{Linear convolution: }\; a \star h = \mathcal{F}^{-1}\!\left[\mathcal{F}[a]\cdot\mathcal{F}[h]\right],
    \label{eq:conv_theorem}
\end{equation}

where the transforms are taken on a zero-padded array of length $M \geq 2N - 1$ (rounded up to the next power of two to allow the FFT) to prevent circular aliasing. The full frFFT algorithm is:

\begin{enumerate}
    \item Compute $a_j = x_j e^{-i\pi\alpha j^2}$ (phase-modulate input by conjugate chirp).
    \item Zero-pad $a$ to length $M = 2^{\lceil\log_2(2N-1)\rceil}$.
    \item Build the chirp kernel $h_{\text{pad}}$ of length $M$ in circular order: $h[0], h[1], \ldots, h[N-1]$ followed by zeros, then $h[N-1], \ldots, h[1]$.
    \item Compute $Y = \texttt{IFFT}(\texttt{FFT}(a_{\text{pad}}) \cdot \texttt{FFT}(h_{\text{pad}}))$.
    \item Multiply the first $N$ outputs: $y_k = e^{-i\pi\alpha k^2} Y_k$, $k = 0, \ldots, N-1$.
\end{enumerate}

This is implemented in the static method \texttt{\_frfft\_batch()} (lines 344--403), where the computation is performed simultaneously for all $J$ trajectories (the \textit{batch} dimension), with \texttt{NumPy}'s vectorised FFT operating along the last axis.

\paragraph{Decoupling in the simulator.} In \texttt{get\_pdf()}, the fractional parameter is set by default to $\alpha = 2/N_\text{fft}$, giving the spatial spacing

\begin{equation}
    \lambda = \frac{2\pi\alpha}{\eta} = \frac{4\pi}{N_\text{fft}\,\eta}.
    \label{eq:lambda_frfft}
\end{equation}

For $N_\text{fft} = 4096$ and $\eta = 0.25$, this yields $\lambda \approx 0.0031$, which is two orders of magnitude finer than the $\lambda_{\text{std}} = 2\pi/(N_\text{fft}\,\eta) \approx 0.0061$ that the standard FFT would enforce. The user can further tighten the grid by choosing a smaller $\alpha$ without any effect on the frequency resolution $\eta$.

% ─────────────────────────────────────────────────────────────────────────────
\section{The PDF Computation Pipeline}
\label{sec:pdf_pipeline}

The method \texttt{get\_pdf()} assembles the above components into a complete PDF computation. The following describes each stage of the pipeline.

\subsection{Grid Initialisation}

For each of the $J$ trajectories, a lower bound $b_j$ for the log-price evaluation grid is determined from an approximate mean and standard deviation of the forecast distribution. Since the exact moments depend on the full path of $v_t$, an approximation using the long-run variance $\theta$ is used:

\begin{align}
    \text{approx\_mean}_j &= X_{0,j} + \left(\mu_j - \tfrac{1}{2}\theta_j\right)\tau, \label{eq:approx_mean}\\
    \text{approx\_std}_j  &= \sqrt{\max(\theta_j, 0)\,\tau}, \label{eq:approx_std}\\
    b_j &= \text{approx\_mean}_j - c\cdot\text{approx\_std}_j, \label{eq:lower_bound}
\end{align}

where $c = 5$ (the parameter \texttt{x\_width}) ensures that the grid spans approximately $\pm 5$ standard deviations around the forecast mean, capturing virtually all of the probability mass.

\subsection{Phase Shift for the Lower Bound}

Equation (\ref{eq:gilpelaez_discrete}) evaluated at $x_k = b + k\lambda$ factors as

\begin{equation}
    f(b + k\lambda) \approx \frac{\eta}{\pi}\,\text{Re}\!\left[\sum_{j=0}^{N_\text{fft}-1} \underbrace{\bigl(w_j\,\varphi(j\eta)\,e^{-ij\eta b}\bigr)}_{\text{frFFT input}} \cdot e^{-i2\pi\alpha jk}\right].
    \label{eq:phase_shifted}
\end{equation}

The factor $e^{-ij\eta b}$ shifts the output grid's origin to $b$ without changing the frequency-domain data. In the code, this is the \texttt{phase} array (line 511), computed as \texttt{np.exp(-1j * u * b)}, which is multiplied element-wise with the weighted CF before passing to \texttt{\_frfft\_batch()}.

\subsection{PDF Extraction and Clipping}

After the frFFT, the approximate PDF values are

\begin{equation}
    f(x_k) \approx \frac{\eta}{\pi}\,\text{Re}\!\left[y_k\right], \quad k = 0, \ldots, N_\text{fft}-1.
    \label{eq:pdf_extract}
\end{equation}

Numerical inversion can produce small negative values due to discretisation error and finite truncation of the frequency integral. These artefacts are removed by clipping: \texttt{pdf\_fine = np.maximum(pdf\_fine, 0.0)}, which has negligible effect on well-resolved distributions.

\subsection{Binning and Normalisation}

The fine PDF grid (of spacing $\lambda$, typically $\sim 0.003$) is resampled onto a coarser shared bin structure with $n_\text{bins}$ equally spaced bins spanning the global range of all $J$ trajectories. Each fine-grid value contributes probability mass $f(x_k)\cdot\lambda$ to the bin containing $x_k$, computed via \texttt{np.histogram} with density weights. The resulting histogram is then normalised row-by-row so that each row sums to one, yielding a proper discrete probability distribution over $n_\text{bins}$ categories.

\subsection{Raw Moment Output}

When \texttt{n\_bins=None}, the method instead computes the mean and standard deviation directly from the fine frFFT grid using the rectangle rule:

\begin{align}
    \mathbb{E}[X] &\approx \frac{\sum_k f(x_k)\,x_k\,\lambda}{\sum_k f(x_k)\,\lambda}, \label{eq:raw_mean}\\
    \text{Var}(X) &\approx \frac{\sum_k f(x_k)\,(x_k - \mathbb{E}[X])^2\,\lambda}{\sum_k f(x_k)\,\lambda}. \label{eq:raw_var}
\end{align}

The output is then a $(J, 2)$ array with columns containing the mean and standard deviation for each trajectory.

% ─────────────────────────────────────────────────────────────────────────────
\section{Monte Carlo Fallback and Validation}
\label{sec:mc}

The method \texttt{\_mc\_pdf()} provides a reference implementation based on direct Monte Carlo simulation. For each of the $J$ starting conditions $(X_{0,j}, v_{0,j})$, a large number \texttt{mc\_sims} of independent forward paths are generated using the full-truncation Euler--Maruyama scheme over \texttt{n\_steps} steps. The terminal log-prices are then histogrammed onto the shared bin structure to produce an empirical approximation of the PDF.

The Monte Carlo estimator is consistent (converges to the true PDF as \texttt{mc\_sims} $\to\infty$) but slow: its computational cost is $\mathcal{O}(J \cdot \texttt{mc\_sims} \cdot \texttt{n\_steps})$. The frFFT method, by contrast, evaluates the exact Heston CF and inverts it in $\mathcal{O}(J \cdot N_\text{fft}\log N_\text{fft})$ operations, making it orders of magnitude faster for large $J$.

The \texttt{plot\_pdf\_comparison()} method overlays the frFFT PDF against a Monte Carlo histogram for visual validation, providing a direct check that the numerical inversion is accurate. Typical agreement between the two methods is excellent for smooth, unimodal distributions and degrades only when the distribution has very heavy tails that require a wider frequency grid or larger $N_\text{fft}$.

% ─────────────────────────────────────────────────────────────────────────────
\section{Parameter Sampling Strategy}
\label{sec:params}

The \texttt{HestonSimulator} supports two modes for specifying each parameter across the $J$ simulated trajectories. Each model parameter (e.g.\ $\mu$, $\kappa$, $\theta$, $\sigma_v$, $\rho$, $v_0$, $X_0$) is specified as either a \texttt{tuple} $(lo, hi)$ or a \texttt{list}.

If a \texttt{tuple} is provided, $J$ values are drawn independently and uniformly from $[\textit{lo},\, \textit{hi}]$:

\begin{equation}
    \theta_j \sim \mathcal{U}[\textit{lo},\, \textit{hi}], \quad j = 1, \ldots, J.
\end{equation}

If a \texttt{list} is provided, the values are used directly. If the list contains a single value, it is broadcast to all $J$ trajectories (constant parameter). This flexibility allows the simulator to generate a diverse dataset of trajectories spanning a wide region of the parameter space in a single call, which is useful for training machine learning models that must generalise across different market regimes.

The random number generator is seeded via \texttt{numpy.random.default\_rng(seed)}, ensuring full reproducibility.

% ─────────────────────────────────────────────────────────────────────────────
\section{Summary}
\label{sec:summary}

The \texttt{HestonSimulator} integrates the following chain of mathematical techniques into a single coherent pipeline. The Heston SDE system (\ref{eq:logprice_sde})--(\ref{eq:correlation}) is discretised using correlated Brownian increments constructed via Cholesky decomposition, with the full-truncation Milstein scheme for the variance and Euler--Maruyama for the log-price. The conditional PDF of future log-prices is obtained by evaluating the closed-form Heston CF in its numerically stable Little Trap formulation and then applying the Gil-Pelaez inversion formula (\ref{eq:gilpelaez_discrete}) with Simpson quadrature weights. The resulting discrete sum is evaluated efficiently across all $N_\text{fft}$ spatial grid points simultaneously by expressing it as a fractional DFT and computing it via Bluestein's chirp convolution algorithm, which reduces the cost from $\mathcal{O}(N_\text{fft}^2)$ to $\mathcal{O}(N_\text{fft}\log N_\text{fft})$ while completely decoupling the frequency resolution $\eta$ from the spatial resolution $\lambda$.